\documentclass[10pt,a4paper]{amsart}
\usepackage{amsmath,amssymb,amsthm,graphicx}
% For commutative diagrams we use Xy-pic
\usepackage[all]{xy}

% We use 2cell for 2-commutative diagrams.
\xyoption{2cell}
% Theorem environments.
%
\theoremstyle{plain}
\newtheorem{theorem}[subsection]{Theorem}
\newtheorem{proposition}[subsection]{Proposition}
\newtheorem{lemma}[subsection]{Lemma}

\theoremstyle{definition}
\newtheorem{definition}[subsection]{Definition}
\newtheorem{example}[subsection]{Example}
\newtheorem{exercise}[subsection]{Exercise}
\newtheorem{situation}[subsection]{Situation}

\theoremstyle{remark}
\newtheorem{remark}[subsection]{Remark}
\newtheorem{remarks}[subsection]{Remarks}

\numberwithin{equation}{subsection}

% Macros
%
\def\lim{\mathop{\mathrm{lim}}\nolimits}
\def\colim{\mathop{\mathrm{colim}}\nolimits}
\def\Spec{\mathop{\mathrm{Spec}}}
\def\Hom{\mathop{\mathrm{Hom}}\nolimits}
\def\Ext{\mathop{\mathrm{Ext}}\nolimits}
\def\SheafHom{\mathop{\mathcal{H}\!\mathit{om}}\nolimits}
\def\SheafExt{\mathop{\mathcal{E}\!\mathit{xt}}\nolimits}
\def\Sch{\mathit{Sch}}
\def\Mor{\mathop{\mathrm{Mor}}\nolimits}
\def\Ob{\mathop{\mathrm{Ob}}\nolimits}
\def\Sh{\mathop{\mathit{Sh}}\nolimits}
\def\NL{\mathop{N\!L}\nolimits}
\def\CH{\mathop{\mathrm{CH}}\nolimits}
\def\proetale{{pro\text{-}\acute{e}tale}}
\def\etale{{\acute{e}tale}}
\def\QCoh{\mathit{QCoh}}
\def\Ker{\mathop{\mathrm{Ker}}}
\def\Im{\mathop{\mathrm{Im}}}
\def\Coker{\mathop{\mathrm{Coker}}}
\def\Coim{\mathop{\mathrm{Coim}}}

% Boxtimes
%
\DeclareMathSymbol{\boxtimes}{\mathbin}{AMSa}{"02}

%
% Macros for moduli stacks/spaces
%
\def\QCohstack{\mathcal{QC}\!\mathit{oh}}
\def\Cohstack{\mathcal{C}\!\mathit{oh}}
\def\Spacesstack{\mathcal{S}\!\mathit{paces}}
\def\Quotfunctor{\mathrm{Quot}}
\def\Hilbfunctor{\mathrm{Hilb}}
\def\Curvesstack{\mathcal{C}\!\mathit{urves}}
\def\Polarizedstack{\mathcal{P}\!\mathit{olarized}}
\def\Complexesstack{\mathcal{C}\!\mathit{omplexes}}
% \Pic is the operator that assigns to X its picard group, usage \Pic(X)
% \Picardstack_{X/B} denotes the Picard stack of X over B
% \Picardfunctor_{X/B} denotes the Picard functor of X over B
\def\Pic{\mathop{\mathrm{Pic}}\nolimits}
\def\Picardstack{\mathcal{P}\!\mathit{ic}}
\def\Picardfunctor{\mathrm{Pic}}
\def\Deformationcategory{\mathcal{D}\!\mathit{ef}}
\usepackage{cleveref}
\crefname{originallemma}{original lemma}{original lemmas}
\Crefname{originallemma}{Original lemma}{Original lemmas}
\makeatletter
\newlabel{sites-modules-lemma-reflect-surjections}{{04DA}{source}{}{sourcepointer-1}{}}
\newlabel{sites-modules-lemma-reflect-surjections@cref}{{[originallemma][0][]04DA}{[1][0][]source}{}{}{}}
\makeatother
\title{Consequences of property (A) for the small pushforward}
\author{The Stacks Project Authors}
\date{stacks-project:etale-cohomology.tex:local-snapshot-20261011; GNU FDL 1.2 or later}
\begin{document}
\maketitle
\noindent\textit{Editorial note.} All delivered mathematical text is the author's own, copied from the pinned Stacks Project chapter file etale-cohomology.tex; no statement or proof step was rewritten, reordered or removed.\par\medskip
\noindent\textit{Editorial note.} The proof is the author's complete argument: reflection of injections and of surjections is proved by the local computation, and the last two clauses are deduced from the first as the author states.\par\medskip
\noindent\textit{Editorial note.} Every labelled reference occurring inside the counted spans is carried in assembly.external\_references as a non-load-bearing navigation pointer (same-chapter ones at their definition site, cross-chapter ones at the invocation inside the counted span), so no load-bearing mathematics is replaced by a stub. No bibliography citation occurs inside the counted spans.\par\medskip
\noindent\textit{Editorial note.} The checked chapter list for same-name collisions is the set of stacks chapters staged in this repository: chow, cohomology, divisors, etale-cohomology, flat, more-algebra, more-morphisms, morphisms, spaces-more-morphisms, varieties together with this card's chapter etale-cohomology.\par\medskip
\noindent\textit{Editorial note.} The preamble is the chapter's own native preamble excerpt already used for this source: the Xy-pic setup (native preamble lines 14-18) and the native theorem declarations with the mathematical macros (native preamble lines 154-217).\par\medskip
\noindent\textit{Allowed standard prerequisites (editorial).}\begin{itemize}
\item Property (A) for a morphism of schemes. Same chapter, the formulation of property (A) delivered verbatim as a context span in this card (source lines 5708-5726): etale locally on X, an etale neighbourhood u of U sits inside the base change of an etale V over Y. Supplies the hypothesis and the local descriptions used throughout the proof.
\item Small pushforward f\_\{small,*\} and reflection of injections and surjections. Same chapter, the functor f\_\{small,*\} on abelian sheaves and the meaning of reflecting injections and surjections. Supplies the properties proved in parts (1) to (3) and the auxiliary result about reflecting surjections cited at the end.
\end{itemize}\medskip
\ifdefined\hypertarget\hypertarget{sourcepointer-1}{}\fi
\noindent\textit{Source reference.} Original lemma 04DA: sources/stacks-etale-cohomology/pinned-original-source.tex, source lines 5791--5791. Navigation pointer invoked at proof line 5791; the counted span contains the pointer only and the cited result is not reproduced inside it, so the pointer is registered as a non-load-bearing navigation stub in the established form for Stacks material. No bibliography citation is involved.\par\medskip
\noindent
Let $f : X \to Y$ be a morphism of schemes.
Here is a list of conditions we will consider in the following:
\begin{enumerate}
\item[(A)] For every \'etale morphism $U \to X$ and $u \in U$ there exist
an \'etale morphism $V \to Y$ and a disjoint union decomposition
$X \times_Y V = W \amalg W'$ and a morphism $h : W \to U$ over $X$
with $u$ in the image of $h$.
\item[(B)] For every $V \to Y$ \'etale, and every \'etale covering
$\{U_i \to X \times_Y V\}$ there exists an \'etale covering
$\{V_j \to V\}$ such that for each $j$ we have
$X \times_Y V_j = \coprod W_{ij}$ where $W_{ij} \to X \times_Y V$
factors through $U_i \to X \times_Y V$ for some $i$.
\item[(C)] For every $U \to X$ \'etale, there exists a $V \to Y$ \'etale
and a surjective morphism $X \times_Y V \to U$ over $X$.
\end{enumerate}
It turns out that each of these properties has meaning in terms of
the behaviour of the functor $f_{small, *}$. We will work this
out in the next few sections.

\begin{lemma}
\label{lemma-property-A-implies}
Let $f : X \to Y$ be a morphism of schemes.
Assume (A).
\begin{enumerate}
\item
$f_{small, *} :
\textit{Ab}(X_\etale)
\to
\textit{Ab}(Y_\etale)$
reflects injections and surjections,
\item $f_{small}^{-1}f_{small, *}\mathcal{F} \to \mathcal{F}$
is surjective for any abelian sheaf $\mathcal{F}$ on $X_\etale$,
\item
$f_{small, *} :
\textit{Ab}(X_\etale)
\to
\textit{Ab}(Y_\etale)$
is faithful.
\end{enumerate}
\end{lemma}

\begin{proof}
Let $\mathcal{F}$ be an abelian sheaf on $X_\etale$.
Let $U$ be an object of $X_\etale$. By assumption we can find a
covering $\{W_i \to U\}$ in $X_\etale$ such that each $W_i$ is
an open and closed subscheme of $X \times_Y V_i$ for some object
$V_i$ of $Y_\etale$. The sheaf condition shows that
$$
\mathcal{F}(U) \subset \prod \mathcal{F}(W_i)
$$
and that $\mathcal{F}(W_i)$ is a direct summand of
$\mathcal{F}(X \times_Y V_i) = f_{small, *}\mathcal{F}(V_i)$.
Hence it is clear that $f_{small, *}$ reflects injections.

\medskip\noindent
Next, suppose that $a : \mathcal{G} \to \mathcal{F}$ is a map of
abelian sheaves such that $f_{small, *}a$ is surjective. Let
$s \in \mathcal{F}(U)$ with $U$ as above. With $W_i$, $V_i$ as
above we see that it suffices to show that $s|_{W_i}$ is \'etale locally
the image of a section of $\mathcal{G}$ under $a$. Since $\mathcal{F}(W_i)$
is a direct summand of $\mathcal{F}(X \times_Y V_i)$
it suffices to show that for any $V \in \Ob(Y_\etale)$
any element $s \in \mathcal{F}(X \times_Y V)$
is \'etale locally on $X \times_Y V$ the image of a section of
$\mathcal{G}$ under $a$. Since
$\mathcal{F}(X \times_Y V) = f_{small, *}\mathcal{F}(V)$
we see by assumption that there exists a covering $\{V_j \to V\}$ such that
$s$ is the image of
$s_j \in f_{small, *}\mathcal{G}(V_j) = \mathcal{G}(X \times_Y V_j)$.
This proves $f_{small, *}$ reflects surjections.

\medskip\noindent
Parts (2), (3) follow formally from part (1), see
Modules on Sites, Lemma \ref{sites-modules-lemma-reflect-surjections}.
\end{proof}

\end{document}
